\chapter*{Wprowadzenie do tomu IV}
\addcontentsline{toc}{chapter}{Wprowadzenie do tomu IV}
W rachunku prawdopodobieństwa ustalaliśmy model i badaliśmy jego konsekwencje. W statystyce kierunek rozumowania jest w pewnym sensie odwrotny: obserwujemy dane i pytamy, jakie właściwości nieznanego mechanizmu ich powstawania możemy na ich podstawie odtworzyć. Nie jest to jednak dosłowne odwracanie funkcji: różne mechanizmy mogą prowadzić do podobnych lub nawet identycznych danych, a przypadkowość oznacza, że na ogół nie otrzymamy pewnego wniosku.

Przykładowo: zaobserwowaliśmy siedem wadliwych produktów w próbie stu sztuk. Liczba $7/100$ jest łatwa do obliczenia, lecz to dopiero początek. Czy produkty wybrano losowo? Czy obserwacje są niezależne? Czy urządzenie pracowało w tym samym trybie? Czy pytamy o całą produkcję, czy o partię z określonej zmiany? Jak zmienia się niepewność wniosku, gdy obserwujemy siedem wad na dziesięć sztuk? W kolejnych rozdziałach będziemy do tego zadania wracać.

\paragraph{Zasada pracy z metodami statystycznymi.}
Dla każdej nowej metody ustalimy: \emph{cel}, \emph{model i założenia}, \emph{wyprowadzenie}, \emph{sposób obliczenia}, \emph{interpretację}, \emph{alternatywy} oraz \emph{sytuacje, w których metoda zawodzi}. Porównywanie metod na tych samych danych nie jest dodatkiem do teorii, tylko częścią jej rozumienia. W tym tomie będziemy rozwijać zarówno wnioskowanie klasyczne, jak i bayesowskie; osobną, systematyczną teorię bayesowską przedstawimy w dalszych częściach.

\paragraph{Co uznajemy za znane.}
Zakładamy pojęcia przestrzeni probabilistycznej (tom I, rozdział~\ref{I-ch:probability}), niezależności i rozkładu (tom I, rozdziały~\ref{I-ch:conditioning} i~\ref{I-ch:random-variables}), wartości oczekiwanej (tom I, rozdział~\ref{I-ch:expectation}), warunkowej wartości oczekiwanej (tom I, rozdział~\ref{I-ch:condexp}) oraz praw wielkich liczb i centralnego twierdzenia granicznego (tom I, rozdziały~\ref{I-ch:lln} i~\ref{I-ch:clt}). Gdy skorzystamy z bardziej zaawansowanego twierdzenia tego tomu (np. Glivenki--Cantellego), podamy dokładnie jego hipotezy. Statystyka nie wymaga przyjmowania nieznanych parametrów za zmienne losowe: w podejściu klasycznym parametr jest ustalony, a losowość przypisujemy obserwacjom.


\section*{Mapa odsyłaczy do pozostałych tomów}
\begin{itemize}
\item \textbf{Tom I — Rachunek prawdopodobieństwa:} rozdział~\ref{I-ch:conditioning} (warunkowanie), rozdział~\ref{I-ch:condexp} (warunkowa wartość oczekiwana), rozdziały~\ref{I-ch:lln}--\ref{I-ch:clt} (twierdzenia graniczne) i rozdział~\ref{I-ch:markov} (łańcuchy Markowa).
\item \textbf{Tom II — Procesy stochastyczne:} rozdział~\ref{II-ch:procesy} (modele trajektorii), rozdział~\ref{II-ch:poisson-process} (proces Poissona), rozdział~\ref{II-ch:wiener} (ruch Browna) oraz rozdział~\ref{II-ch:quadratic-cont} (wariacja kwadratowa).
\item \textbf{Tom III — Analiza stochastyczna:} rozdział~\ref{III-ch:sde-existence} (SDE), rozdział~\ref{III-ch:sde-properties} (aproksymacja numeryczna), rozdział~\ref{III-ch:levy} (procesy Lévy'ego) i rozdział~\ref{III-ch:path-integrals} (funkcjonały trajektorii).
\end{itemize}
Odsyłacze w tej sekcji oraz w akapitach \emph{Powiązania i odsyłacze} odwołują się do pozostałych tomów. W oddzielnych plikach PDF są zapisane jako odsyłacze \texttt{GoToR}, których obsługa zależy od czytnika i jego ustawień bezpieczeństwa. Aby mieć pewność działania wszystkich odsyłaczy także w przeglądarce internetowej, należy używać wydania zbiorczego \texttt{Cztery\_tomy\_polaczone\_aktywne\_linki.pdf}, gdzie są one zwykłymi odsyłaczami wewnętrznymi. Wydanie zbiorcze przygotowuje dostarczony skrypt \texttt{polacz\_tomy.py}. Odsyłacze lokalne w pojedynczym tomie nie wymagają innych plików.

